\section{Table File Format}\label{sec:table-file-format}

Table files are text files that can be created with a text editor or word processor.
On UNIX, the standard \texttt{vi} editor can be used; on a Silicon Graphics
workstation, the \texttt{jot} editor can be used. Table files can also be created with
other software, such as Aero\taosref{20} and Rocket.\taosref{21}

\taossubhead{Values}

Table data is read line by line and is not column dependent. Information must be
given in the correct order so TAOS knows how to interpret it. Lines can be any
number of characters in length, and blank lines are ignored. Names and values in the
file can have up to 72 characters, so they are essentially unlimited in length for the
intended use.

Names and values are separated with white-space characters (such as space, tab, or
carriage return), commas, equal signs, parentheses, colons, greater-than signs, and
less-than signs. Because these special characters are delimiters, they cannot be
embedded within a name or value. For example, the value \texttt{30,000} is not
correct because the comma separates values; TAOS interprets it as two values,
\texttt{30} and \texttt{0}. The correct input is \texttt{30000}.

Input values are free-field rather than column dependent, so indentation can be used
to make a table file easier to understand. The examples in this manual use
indentation to show relationships among values. This style is encouraged because it
reduces input errors.

Values can be entered with or without a decimal point. All values are interpreted as
real numbers; no distinction is made for integers. Values can also be entered in
scientific notation, or e-format. For example, the number 2000 can be entered as
\texttt{2000}, \texttt{2000.0}, or \texttt{2.0e3}.

TAOS converts all characters to lower case before interpreting them as names or
values, so upper- and lower-case letters are equivalent. Lower-case characters are
recommended for readability.

\taossubhead{Comments}

Comments start with the special character \texttt{\#} and continue to the end of the
line. They can be placed anywhere in the file except within a name or value. The
following lines show valid comments:

\begin{lstlisting}[style=taosmanual]
# the entire line is used for comments
0.2 0.3 0.4 0.5 # Mach numbers
\end{lstlisting}

In the first line, the \texttt{\#} symbol is in the first column. Everything after it is
treated as a comment and ignored by TAOS, so the entire line is a comment. In the
second example, the comment begins after the four values; everything from the
\texttt{\#} symbol to the end of the line is ignored.

\taossubhead{State Variables}

As TAOS computes the trajectory of a vehicle, it keeps track of many values. These
values together define the vehicle's state as it moves along the trajectory, so they are
called \emph{state variables}. As noted previously, a table defines a single value: it
contains enough information for TAOS to compute that value. Usually the value is a
function of one or more state variables. For example, thrust is often a function of
time. At each instant along the trajectory, time is known; therefore, the thrust value
can be determined. \Cref{tab:taos-state-variables} lists the state variables that can be
referenced in TAOS tables.

\begingroup
\small
\setlength{\LTpre}{0.8\baselineskip}
\setlength{\LTpost}{0.8\baselineskip}
\renewcommand{\arraystretch}{1.03}
\begin{longtable}{@{}>{\ttfamily}p{0.105\textwidth}p{0.31\textwidth}>{\ttfamily}p{0.105\textwidth}p{0.31\textwidth}@{}}
\caption{TAOS state variables.}\label{tab:taos-state-variables}\\
\toprule
\normalfont Variable & Description & \normalfont Variable & Description \\
\midrule
\endfirsthead
\caption[]{TAOS state variables (continued).}\\
\toprule
\normalfont Variable & Description & \normalfont Variable & Description \\
\midrule
\endhead
\bottomrule
\endfoot
alpha   & Angle of attack & pitchgd & Geodetic pitch angle \\
alphat  & Total angle of attack & pitchi & Inertial pitch angle \\
alt\limitedstate & Geodetic altitude & power & Power setting \\
altdt\limitedstate & Geodetic altitude rate & pres\limitedstate & Atmospheric pressure \\
bankgc & Geocentric bank angle & psigc\limitedstate & Geocentric heading angle \\
bankgd & Geodetic bank angle & psigd\limitedstate & Geodetic heading angle \\
beta & Sideslip angle & range & Range \\
betae & Euler sideslip angle & rcm\limitedstate & Distance from earth center \\
cg & Center of gravity & rcmdt\limitedstate & Rate from earth center \\
dynprs\limitedstate & Dynamic pressure & reypft & Reynolds number per foot \\
ep1 & Thrust-vector angle 1 & rho\limitedstate & Atmospheric density \\
ep2 & Thrust-vector angle 2 & rollgc & Geocentric roll angle \\
gamgc\limitedstate & Geocentric flight-path angle & rollgd & Geodetic roll angle \\
gamgd\limitedstate & Geodetic flight-path angle & rolli & Inertial roll angle \\
grmark & Ground range since last reset & segment & Segment number \\
grseg & Segment ground range & sndspd\limitedstate & Atmospheric speed of sound \\
latgc\limitedstate & Geocentric latitude & sref & Aerodynamic reference area \\
latgcdt & Geocentric latitude rate & temp\limitedstate & Atmospheric temperature \\
latgd\limitedstate & Geodetic latitude & thrust & Total thrust \\
latgddt & Geodetic latitude rate & time\limitedstate & Time \\
long\limitedstate & Longitude & tmark\limitedstate & Time since last reset \\
longdt & Longitude rate & tseg\limitedstate & Segment time \\
mach & Mach number & vair & Velocity relative to air \\
mass\limitedstate & Vehicle mass & vel\limitedstate & Inertial velocity \\
nu\limitedstate & Atmospheric kinematic viscosity & vgr & Velocity relative to ground \\
plength & Path length & visc\limitedstate & Atmospheric viscosity \\
plmark & Path length since last reset & wt\limitedstate & Vehicle weight \\
plseg & Segment path length & yawgc & Geocentric yaw angle \\
phi & Windward meridian & yawgd & Geodetic yaw angle \\
pitchgc & Geocentric pitch angle & yawi & Inertial yaw angle \\
\end{longtable}
\endgroup

Wind and center-of-gravity tables can only be functions of the state variables marked
with a dagger (\limitedstate) in \cref{tab:taos-state-variables}. These tables are
evaluated before all state variables are known. Tables also cannot be functions of
themselves; for example, a thrust table cannot be a function of \statevar{thrust}.

The default units for state variables are given in the appendix. They can be changed
in the problem file with the \problemblock{*units/fmt} data block
(\cref{sec:block-units-fmt}).

\taossubhead{User-Defined Variables}

There are times when the state variables alone do not provide enough flexibility. For
example, there may be three thrust-versus-time curves for three versions of a rocket
motor. The curves can be placed in separate tables, or they can be placed in one
table. If they are placed in one table, thrust becomes a function of both time and
another variable that might be called \texttt{version}. This new variable is a
\emph{user-defined variable}: any name can be used except those reserved for the
state variables in \cref{tab:taos-state-variables}.

When TAOS detects a user-defined variable, its value must be given in the problem
file. For example,

\begin{lstlisting}[style=taosmanual]
*prop thrust=(stage1) version=2
\end{lstlisting}

says to use a table called \tableid{stage1} for thrust and to set the user-defined
variable \texttt{version} equal to two. The value of \texttt{version} is constant within
a trajectory segment, but it can be surveyed like any other numerical input value
(\cref{sec:block-survey}).
